\documentclass[11pt]{amsart}
\usepackage[T1]{fontenc}
\usepackage{lmodern}
\usepackage{amsmath,amssymb,mathtools}
\usepackage[margin=1.1in]{geometry}
\usepackage{microtype}
\usepackage[colorlinks=true,linkcolor=blue,citecolor=blue,urlcolor=blue]{hyperref}
\hypersetup{pdftitle={Moment exponents for Gaussian directed polymers in three dimensions},pdfauthor={Kaiwen Jin}}
\newcommand{\E}{\mathbb E}
\newcommand{\Pp}{\mathbb P}
\newcommand{\R}{\mathbb R}
\newcommand{\Z}{\mathbb Z}
\newcommand{\ps}{p^*}
\newtheorem{theorem}{Theorem}[section]
\newtheorem{proposition}[theorem]{Proposition}
\newtheorem{lemma}[theorem]{Lemma}
\theoremstyle{remark}
\newtheorem{remark}[theorem]{Remark}
\numberwithin{equation}{section}
\title{Moment exponents for Gaussian directed polymers in three dimensions}
\author{Kaiwen Jin}
\email{yui.ui.siki@gmail.com}
\date{October 15, 2026}
\keywords{Directed polymers, Gaussian disorder, moment exponent, weak disorder}
\begin{document}
\begin{abstract}
For the directed polymer based on simple random walk in $\Z^3$ with iid
standard Gaussian disorder, let $p^*(\beta)$ be the supremum of the orders for
which the normalized partition function has moments bounded uniformly in the
polymer length. We prove that $\beta\mapsto p^*(\beta)$ is a continuous strictly
decreasing bijection from $(0,\beta_c]$ onto $[5/3,\infty)$, where $\beta_c$
is the weak-disorder critical point. The main step is local stability of every
strictly subcritical moment above order two under an increase of $\beta$.
A concave power of a terminal polymer moment has a one-step deficit proportional
to the endpoint overlap; Gaussian differentiation bounds the perturbation by
the same overlap. This yields a moment bound for all lengths. Gaussian
hypercontractivity and the established critical and one-sided continuity
results then give the full statement.
\end{abstract}
\maketitle

\section{Introduction}
A directed polymer in a random environment weights random-walk trajectories
according to the disorder they encounter. Its normalized partition function
is a positive mean-one martingale. The weak-disorder phase is characterized by
a positive martingale limit, but its fluctuations need not have a bounded
second moment. The threshold for higher moments gives a finer description of
this phase; see \cite[Section~3.3.2]{Berger} for the moment-spectrum question
and its relation to the hierarchy of $L^p$ critical points.

We consider simple random walk $S=(S_n)_{n\ge0}$ from zero on $\Z^3$, with
each of its six nearest-neighbor steps having probability $1/6$. Let
$(\omega_{n,x})_{n\ge1,x\in\Z^3}$ be iid standard real Gaussians independent
of the walk. For inverse temperature $\beta>0$, define
\begin{equation}\label{eq:partition}
 W_N^\beta=E_0\exp\!\left\{\sum_{n=1}^N
 (\beta\omega_{n,S_n}-\beta^2/2)\right\},\qquad N\ge0,
\end{equation}
where $E_0$ denotes walk expectation and $\E$ denotes disorder expectation.
We set $W_0^\beta=1$ and write $W_\infty^\beta$ for the almost sure martingale
limit. The moment exponent and the weak-disorder critical point are
\begin{align}
 \ps(\beta)&=\sup\{u\ge1:\sup_{N\ge0}\E[(W_N^\beta)^u]<\infty\},
 \label{eq:pstar}\\
 \beta_c&=\sup\{\beta>0:\Pp(W_\infty^\beta>0)=1\}.
 \label{eq:betac}
\end{align}
All moment orders are real. The supremum in \eqref{eq:pstar} does not assert
boundedness of the moment at its endpoint.

Several ingredients of the spectrum are known for much more general
disorder. Weak disorder holds below the critical point and also at criticality
for the simple walk in dimension at least three; in dimension three the
critical exponent is $5/3$ \cite[Theorem~2.1 and Corollary~2.5]{Critical}.
Junk and Lacoin prove left continuity wherever weak disorder holds
\cite[Corollary~2.11(i)]{Tail}. Junk proves right continuity when
$5/3<\ps(\beta)\le2$ \cite[Theorem~1.2(i)]{Junk}. The strict-monotonicity
statement in part (ii) of that theorem assumes finite-support disorder and
therefore does not apply to the Gaussian law. Gaussian hypercontractivity
provides a direct comparison between different inverse temperatures, as we
recall in Section~\ref{sec:comparison}. To obtain continuity of the whole
Gaussian spectrum we prove upward stability of moment bounds in the range
$\ps(\beta)>2$.

\begin{theorem}\label{thm:main}
For the Gaussian model \eqref{eq:partition}, $0<\beta_c<\infty$, and
\[
 (0,\beta_c]\ni\beta\longmapsto\ps(\beta)\in[5/3,\infty)
\]
is a continuous strictly decreasing bijection. Continuity at $\beta_c$ is
from the left.
\end{theorem}

This proves the Gaussian three-dimensional instance of the continuity and
strict-monotonicity conjecture stated in \cite[Conjecture~3.32]{Berger}.
It describes every finite real value of the moment exponent on the original
weak-disorder interval. It does not settle boundedness at the threshold
moment order.

The stability argument uses the limiting partition functions started at
different sites in a common future environment. Their joint law gives a
harmonic identity for the terminal $p$ moment of any initial probability
vector. Taking its power $s/p<1$ creates a concavity deficit. Positive
association and a Gaussian change of measure bound that deficit from below
by a constant times the squared $\ell^2$ norm of the predicted endpoint law.
The derivative with respect to the disorder variance has the same factor.
A small upward change of variance is therefore absorbed uniformly, even for
very spread-out endpoint laws. Iterating the resulting inequality gives the
$s$ moment bound for every polymer length. The strict gaps
$2<s<p<q<\ps(\beta)$ provide both the differentiability and the density
integrability needed in this argument.

Section~\ref{sec:comparison} records the precise existing continuity results
and proves the Gaussian comparison. Section~\ref{sec:stability} establishes
local moment stability and completes the proof of Theorem~\ref{thm:main}.
Throughout, expectations without a walk subscript refer to the disorder;
additional subscripts indicate which independent variables are averaged.
\section{Gaussian comparison and continuity facts}\label{sec:comparison}

Write $W_N^{\beta,x}$ for the partition function started at $x$, and put
$W_0^{\beta,x}=1$. With respect to the disorder filtration
$\mathcal F_N=\sigma(\omega_{n,x}:n\le N)$, these are positive mean-one
martingales. Indeed, averaging the fresh disorder slice gives one. For every
real $u\ge1$, Jensen's inequality over the walk paths gives
\begin{equation}\label{eq:finite-moments}
 \E[(W_N^{\beta,x})^u]\le \exp\{N(u^2-u)\beta^2/2\}.
\end{equation}
The nonnegative martingale convergence theorem therefore defines
$W_\infty^{\beta,x}$ almost surely; see, for example,
\cite[Theorem~4.2.12]{Durrett}. We will repeatedly use the elementary consequence
of the definition of $\ps$: every $1\le u<\ps(\beta)$ is a uniformly bounded
moment order. To see this, choose a bounded order $v>u$ and apply H\"older's
inequality. No assertion at $u=\ps(\beta)$ is needed.

There is a positive weak-disorder interval and a parameter $\beta_\ell>0$ at
which the martingale is bounded in $L^2$ \cite[Theorem~A]{Junk}. In particular,
$\beta_c>0$. The opposite bound is elementary in the Gaussian case. For
$0<\theta<1$, summing the $6^N$ path contributions and using subadditivity yields
\[
 \E[(W_N^\beta)^\theta]
 \le \exp\{N(1-\theta)(\log6-\theta\beta^2/2)\}.
\]
If $\beta^2>2\log6$, choose $\theta<1$ with $\theta\beta^2/2>\log6$.
Fatou's lemma then gives $W_\infty^\beta=0$ almost surely. Thus
\begin{equation}\label{eq:critical-finite}
 0<\beta_c\le\sqrt{2\log6}<\infty.
\end{equation}
Similarly, for every $u>1$, the pointwise inequality
$(\sum_i a_i)^u\ge\sum_i a_i^u$ gives
\begin{equation}\label{eq:diagonal}
 \E[(W_N^\beta)^u]
 \ge \exp\{N(u-1)(u\beta^2/2-\log6)\}.
\end{equation}
Hence $\ps(\beta)$ is finite for every $\beta>0$; if $\ps(\beta)>1$, it is at
most $2\log6/\beta^2$.

Let $\beta_2=\sup\{\beta>0:\sup_N\E[(W_N^\beta)^2]<\infty\}$ denote
the $L^2$ critical point. We use the following established results, with their
original disorder assumptions specialized to standard Gaussians:
\begin{enumerate}
\item Weak disorder holds for $0<\beta\le\beta_c$, and
\begin{equation}\label{eq:critical-value}
 \ps(\beta_c)=\frac53.
\end{equation}
The phase intervals are stated in \cite[Theorem~A]{Junk}. Weak disorder at
criticality follows from \cite[Theorem~2.1 and the discussion after
Remark~2.2]{Critical}; its simple-walk discussion includes the strict inequality
$\beta_c>\beta_2$ in dimension three. The value in
\eqref{eq:critical-value} is \cite[Corollary~2.5]{Critical}.
\item At every parameter where weak disorder holds,
\begin{equation}\label{eq:left-continuity}
 \lim_{b\uparrow\beta}\ps(b)=\ps(\beta)
\end{equation}
by \cite[Corollary~2.11(i)]{Tail}.
\item If $5/3<\ps(\beta)\le2$, then $\ps$ is right continuous at $\beta$
by \cite[Theorem~1.2(i)]{Junk}.
\end{enumerate}
The Gaussian law has all real exponential moments, so the unbounded-disorder
results cited here apply. In \cite{Tail}, the threshold is written as the
infimum of the orders whose uniform moments are infinite. This is the same
number as \eqref{eq:pstar}: H\"older makes the bounded orders a downward interval,
and \eqref{eq:diagonal} supplies finite orders outside it. Neither convention
determines whether its endpoint belongs to that interval.

Gaussian hypercontractivity supplies a quantitative comparison between
temperatures. For a function of a finite standard Gaussian vector, define
\[
 T_\rho f(z)=\E_G f(\rho z+\sqrt{1-\rho^2}\,G),\qquad 0\le\rho\le1,
\]
where $G$ is an independent standard Gaussian vector of the same dimension.
Nelson's inequality states
\begin{equation}\label{eq:hypercontractivity}
 \|T_\rho f\|_v\le\|f\|_u,
 \qquad 1<u\le v<\infty,\quad
 \rho\le\sqrt{\frac{u-1}{v-1}}.
\end{equation}
See \cite[Theorem~1.1 in the author-generated copy]{Kemp}. In Ornstein--Uhlenbeck time $a$, this is the
semigroup with $\rho=e^{-a}$; the Mehler formula and the corresponding exponent
$v=1+(u-1)e^{2a}$ are given in \cite[Lemma~1.8 and Theorem~2.1]{ChafaiLehec}.
The constant is independent of the dimension.

\begin{proposition}\label{prop:comparison}
For $0<b<\beta$ with $\ps(\beta)>1$,
\begin{equation}\label{eq:threshold-comparison}
 \ps(b)\ge1+\frac{\beta^2}{b^2}\bigl(\ps(\beta)-1\bigr).
\end{equation}
Consequently, $\ps$ is strictly decreasing on $(0,\beta_c]$ and
$\lim_{b\downarrow0}\ps(b)=\infty$.
\end{proposition}

\begin{proof}
For fixed $N$, only finitely many disorder coordinates enter $W_N^\beta$.
Every path uses $N$ distinct coordinates, since their time indices differ.
Integrating the added Gaussian vector in $T_\rho$ therefore gives
\begin{equation}\label{eq:mehler-polymer}
 T_\rho W_N^\beta=W_N^{\rho\beta}.
\end{equation}
More explicitly, the integration contributes
$\exp\{N\beta^2(1-\rho^2)/2\}$, changing the compensator from
$N\beta^2/2$ to $N(\rho\beta)^2/2$ in each path term. The path sum is finite.
Jensen's inequality for $T_\rho$, followed by Gaussian invariance, first shows
that $\ps$ is nonincreasing.

Now take $1<u<\ps(\beta)$ and set $\rho=b/\beta$,
$v=1+(\beta/b)^2(u-1)$. Equations \eqref{eq:hypercontractivity} and
\eqref{eq:mehler-polymer} transfer the uniform $u$ moment at $\beta$ to a
uniform $v$ moment at $b$. Letting $u\uparrow\ps(\beta)$ proves
\eqref{eq:threshold-comparison}. This passes to a supremum of admissible
orders, rather than asserting a bounded moment at that supremum.

Monotonicity and \eqref{eq:critical-value} give
$\ps(\beta)\ge5/3$ for $\beta\le\beta_c$. Since $\beta^2/b^2>1$,
\eqref{eq:threshold-comparison} is then strictly larger than $\ps(\beta)$.
Finally, the actually bounded second moment at $\beta_\ell$ and
\eqref{eq:hypercontractivity} give
$\ps(b)\ge1+\beta_\ell^2/b^2$ for $0<b<\beta_\ell$.
\end{proof}

\section{Stability of moments above two}\label{sec:stability}

The remaining step is an upward perturbation of the inverse temperature.
The moment loss in the statement below leaves room for strict concavity and
an integrable change of measure.

\begin{proposition}\label{prop:stability}
Let $\beta_0>0$ and $2<s<\ps(\beta_0)$. There are $\varepsilon>0$ and
$C<\infty$ such that
\begin{equation}\label{eq:stability}
 \sup_{\beta\in[\beta_0,\beta_0+\varepsilon]}
 \sup_{N\ge0}\E[(W_N^\beta)^s]\le C.
\end{equation}
\end{proposition}

\subsection{The common terminal field}
Fix real exponents
\begin{equation}\label{eq:exponents}
 2<s<p<q<\ps(\beta_0),\qquad a=\frac{s}{p}\in(0,1),\qquad r=\frac qp>1.
\end{equation}
For $j\ge0$ and $x\in\Z^3$, let $W_N^{(j),x}$ denote the
$\beta_0$ partition function using disorder times $j+1,\ldots,j+N$ and
started at $x$. These martingales are uniformly bounded in $L^q$, and hence
converge almost surely and in $L^q$ \cite[Theorem~4.4.6]{Durrett}.
Write their limits as $X_x^{(j)}$. On a common probability-one set for the
countable family $(j,x)$, the six-term recursion passes to the limit:
\begin{equation}\label{eq:terminal-recursion}
 X_x^{(j)}=\sum_y K(x,y)\eta_{j+1,y}^{\beta_0}X_y^{(j+1)},
 \qquad \eta_{n,y}^{\beta}=e^{\beta\omega_{n,y}-\beta^2/2},
\end{equation}
where $K(x,y)=1/6$ for nearest neighbors and zero otherwise.
The whole field $X^{(j+1)}$ is independent of the first $j+1$ disorder
slices and has the same joint law as $X^{(0)}$. Its coordinates have mean one.
Put
\[
 C_p=\E[(X_0^{(0)})^p],\qquad C_q=\E[(X_0^{(0)})^q].
\]
These finite marginal moments are unchanged by spatial translation.

We also need positive association. For bounded increasing functions $f,g$
of one random variable $U$, with $U'$ an independent copy,
\[
 \operatorname{Cov}(f(U),g(U))
 =\tfrac12\E[(f(U)-f(U'))(g(U)-g(U'))]\ge0.
\]
Conditioning on one coordinate and using induction proves the same inequality
for increasing functions of finitely many independent coordinates. For a
countable product, condition bounded $f,g$ on the first $m$ coordinates.
Their conditional averages are increasing on the finite product and converge
almost surely to $f,g$ by the increasing-sigma-field convergence theorem
\cite[Theorem~4.6.8]{Durrett}. Bounded convergence passes the covariance
inequality to the limit. Truncation and monotone convergence extend it to
nonnegative functions with finite product expectation.

Each $X_x^{(j)}$ may be represented by the extended nonnegative liminf of its
finite partition functions. This is increasing in every future disorder
coordinate and agrees almost surely with the $L^q$ limit. Thus the joint
terminal field is associated. This construction uses its common future;
its spatial coordinates are not replaced by independent copies.

Let $\mu$ be any finitely supported probability vector on $\Z^3$. To measure
its terminal moment, define, using a copy $X$ of the whole field $X^{(0)}$,
\begin{equation}\label{eq:h-H}
 h(\mu)=\E_X\Big(\sum_x\mu_xX_x\Big)^p,
 \qquad H(\mu)=h(\mu)^a.
\end{equation}
Jensen and mean one give the support-independent bounds
\begin{equation}\label{eq:terminal-bounds}
 1\le h(\mu)\le C_p,\qquad 1\le H(\mu)\le C_p^a,
 \qquad \E_X\Big(\sum_x\mu_xX_x\Big)^q\le C_q.
\end{equation}
On any fixed finite-support simplex, $h$ and $H$ are continuous by dominated
convergence: the $p$th power of the sum of the terminal coordinates in that
fixed support is an integrable dominator.
In particular they are measurable at the finite endpoint distributions used
below.

Predict the next endpoint law by $\alpha=\mu K$ and write
\begin{equation}\label{eq:overlap}
 I=\sum_x\alpha_x^2>0,\qquad t_0=\beta_0^2.
\end{equation}
The positive support $J=\{x:\alpha_x>0\}$ is finite. For independent standard
Gaussians $(g_x)_{x\in J}$, independent of $X$, put
\begin{equation}\label{eq:one-step}
 \eta_x(t)=e^{\sqrt t\,g_x-t/2},\qquad
 A_t=\sum_x\alpha_x\eta_x(t),\qquad
 \nu_t(x)=\frac{\alpha_x\eta_x(t)}{A_t}.
\end{equation}
Here $t$ is the variance parameter $\beta^2$, not the semigroup time used in
Section~\ref{sec:comparison}.
For a positive vector $b=(b_x)_{x\in J}$ set
\begin{equation}\label{eq:L-F}
 Z_b=\sum_xb_xX_x,\qquad B=\sum_xb_x,\qquad
 L(b)=\E_X Z_b^p,\qquad F(b)=L(b)^a.
\end{equation}
The joint recursion \eqref{eq:terminal-recursion} and homogeneity imply
\begin{align}
 \E_g L(\alpha\eta(t_0))&=h(\mu),\label{eq:harmonic}\\
 \E_{g,X}\Big(\sum_x\alpha_x\eta_x(t_0)X_x\Big)^q&\le C_q,\label{eq:q-harmonic}\\
 F(\alpha\eta(t))&=A_t^sH(\nu_t).\label{eq:homogeneity}
\end{align}
Indeed, the sum in the first two lines has the law of
$\sum_x\mu_xX_x^{(0)}$. This is where retaining the joint terminal field
matters.

\subsection{Gaussian differentiation}
We first control the variation of $\E_gF(\alpha\eta(t))$ with $t$.
For positive $b$, Jensen gives $B^p\le L(b)\le C_p B^p$.
Differentiation under the $X$ expectation gives
\[
 L_x=p\E_X[X_xZ_b^{p-1}],\qquad
 L_{xx}=p(p-1)\E_X[X_x^2Z_b^{p-2}].
\]
On compact subsets of the positive orthant, the integrand and these
derivatives are bounded by an integrable multiple of
$(\sum_{x\in J}X_x)^p$. Thus the differentiations are valid also for
noninteger $p>2$. H\"older gives
\[
 \E_X[X_xZ_b^{p-1}]\le C_p B^{p-1},\qquad
 \E_X[X_x^2Z_b^{p-2}]\le C_p B^{p-2},
\]
and Cauchy--Schwarz gives
$L_x^2\le p^2L\E_X[X_x^2Z_b^{p-2}]$.
Since $a<1$,
\begin{equation}\label{eq:Fxx}
 F_{xx}=aL^{a-1}L_{xx}-a(1-a)L^{a-2}L_x^2
 \le s(p-1)C_p B^{s-2}.
\end{equation}
The same estimates also give
\begin{equation}\label{eq:absolute-derivatives}
 |F_x|\le s C_p B^{s-1},\qquad
 |F_{xx}|\le\{s(p-1)+a(1-a)p^2\}C_p B^{s-2}.
\end{equation}
Both $p-2$ and $s-2$ are positive.

Finite-dimensional Gaussian integration by parts now yields
\begin{equation}\label{eq:variance-derivative}
 \frac d{dt}\E_gF(\alpha\eta(t))
 =\frac12\E_g\sum_x(\alpha_x\eta_x(t))^2F_{xx}(\alpha\eta(t)).
\end{equation}
For completeness, differentiating before integration by parts gives
$\sum_x b_xF_x(g_x/(2\sqrt t)-1/2)$, with $b=\alpha\eta(t)$.
The identity $\E_g(g_x f)=\E_g(\partial_{g_x}f)$ replaces
$g_xb_xF_x$ by $\sqrt t(b_xF_x+b_x^2F_{xx})$; the terms $b_xF_x$
cancel. On a compact interval of positive $t$, \eqref{eq:absolute-derivatives}
bounds the relevant functions and derivatives by finite Gaussian exponentials
with polynomial factors in $g$. These are integrable and dominate the
difference quotients. They also make the boundary terms vanish against the
Gaussian density. Dominated convergence and Fubini justify the derivative and
integration by parts for each finite $J$.

By H\"older with exponents $s/2$ and $s/(s-2)$,
\[
 \E_g[\eta_x(t)^2 A_t^{s-2}]
 \le (\E\eta_x(t)^s)^{2/s}(\E A_t^s)^{(s-2)/s}
 \le e^{s(s-1)t/2}.
\]
The last inequality follows from Jensen over the probability vector $\alpha$.
Consequently, \eqref{eq:Fxx}--\eqref{eq:variance-derivative} give
\begin{equation}\label{eq:perturbation-derivative}
 \frac d{dt}\E_gF(\alpha\eta(t))
 \le \frac{s(p-1)C_p}{2}e^{s(s-1)t/2} I.
\end{equation}
This bound has the same overlap $I$ as the concavity estimate proved next.

\subsection{A uniform concavity estimate}
Recall $H(\mu)=h(\mu)^a$ and $F(b)=L(b)^a$.
We show that $H(\mu)$ exceeds $\E_gF(\alpha\eta(t_0))$ by a fixed
positive multiple of $I$, uniformly over all finite supports.
Set
\[
 Y=L(\alpha\eta(t_0)),\qquad D=\frac{Y}{h(\mu)}.
\]
By \eqref{eq:harmonic}, $D$ is a positive density of mean one. Since $r=q/p>1$,
conditional Jensen and \eqref{eq:q-harmonic} give
\begin{equation}\label{eq:density-moment}
 \E_gD^r\le\E_gY^r
 \le\E_{g,X}\Big(\sum_x\alpha_x\eta_x(t_0)X_x\Big)^q\le C_q.
\end{equation}
The scalar variable
\[
 G_\alpha=I^{-1/2}\sum_x\alpha_xg_x
\]
is standard normal for every $\alpha$. Let $r'=r/(r-1)$ and define
\begin{equation}\label{eq:M}
 M=C_q^{1/r}(\E|G|^{2r'})^{1/r'},\qquad G\sim N(0,1).
\end{equation}
H\"older gives $\E_g[D G_\alpha^2]\le M$, with no dependence on $J$.

Association of the terminal field and Jensen, using $p-1>1$, imply
\begin{equation}\label{eq:association-lower}
 \E_X[X_xZ_b^{p-1}]
 \ge \E_XX_x\,\E_XZ_b^{p-1}\ge B^{p-1}.
\end{equation}
The product is integrable by the $p$-moment estimates above. Differentiating
$Y$ and integrating by parts in $g_x$ therefore give
\begin{align*}
 \E_g[G_\alpha Y]
 &=\frac{p\beta_0}{\sqrt I}\sum_x\alpha_x^2
   \E_g\!\left[\eta_x(t_0)\E_X(X_x Z_{\alpha\eta(t_0)}^{p-1})\right]\\
 &\ge p\beta_0\sqrt I.
\end{align*}
For the last inequality, \eqref{eq:association-lower} bounds the inner terminal
expectation by $A_{t_0}^{p-1}$. Association in the fresh Gaussian slice and
Jensen then yield
$\E_g[\eta_x(t_0)A_{t_0}^{p-1}]\ge\E_g\eta_x(t_0)\,\E_gA_{t_0}^{p-1}\ge1$.
The growth estimates used in \eqref{eq:variance-derivative} also justify this
integration by parts. Since $h(\mu)\le C_p$, we obtain
\begin{equation}\label{eq:shift}
 \E_g[D G_\alpha]\ge\frac{p\beta_0}{C_p}\sqrt I.
\end{equation}

The original Gaussian mean of $G_\alpha$ is zero. Thus
\[
 \E_g[D G_\alpha]
 =\E_g[(\sqrt D-1)G_\alpha(\sqrt D+1)].
\]
Cauchy--Schwarz, $(\sqrt D+1)^2\le2(D+1)$, and \eqref{eq:M} imply
\begin{equation}\label{eq:hellinger}
 \E_g(\sqrt D-1)^2
 \ge\frac{p^2\beta_0^2}{2C_p^2(1+M)}I.
\end{equation}
Only the $r$ moment of $D$ is used; it need not have a second moment.

For $0<a<1$ and $z\ge0$,
\begin{equation}\label{eq:scalar-concavity}
 1-a+az-z^a\ge\min(a,1-a)(\sqrt z-1)^2.
\end{equation}
If $a\le1/2$, apply the tangent-line inequality to $u^{2a}$ at $u=1$.
If $a\ge1/2$, use the first case with $1-a$ and $1/z$, and multiply by $z$;
the value at zero follows by continuity. Integrate
\eqref{eq:scalar-concavity} at $D$ and use $\E_gD=1$.
By \eqref{eq:harmonic}, $Y=h(\mu)D$, while $h(\mu)^a\ge1$.
Together with \eqref{eq:hellinger}, this gives
\begin{equation}\label{eq:deficit}
 H(\mu)-\E_gF(\alpha\eta(t_0))\ge cI,\qquad
 c=\frac{\min(a,1-a)p^2\beta_0^2}{2C_p^2(1+M)}>0.
\end{equation}

\subsection{Perturbation and all-length iteration}
Keep $\beta_0,s,p,q$ and the resulting $H$ fixed. Define
\[
 Q=\frac{s(p-1)C_p}{2}e^{s(s-1)(t_0+1)/2},\qquad
 \tau=\min\{1,c/(2Q)\}>0.
\]
Integrating \eqref{eq:perturbation-derivative} and using
\eqref{eq:homogeneity} and \eqref{eq:deficit}, for every
$t\in[t_0,t_0+\tau]$ we get
\begin{equation}\label{eq:lyapunov-step}
 \E_g[A_t^sH(\nu_t)]
 \le H(\mu)-cI+Q(t-t_0)I\le H(\mu).
\end{equation}
The constants $c,Q,\tau$ are independent of $\mu$ and its support.
Both the loss and the perturbation contain exactly $I=\sum_x\alpha_x^2$,
so no lower bound on this overlap is required.

Fix $\beta$ with $\beta^2\in[t_0,t_0+\tau]$, and let
\[
 u_n(x)=E_0\!\left[\prod_{j=1}^n\eta_{j,S_j}^\beta;\ S_n=x\right],
 \qquad \mu_n(x)=\frac{u_n(x)}{W_n^\beta}.
\]
These are the actual endpoint weights and endpoint law at the perturbed
parameter. We have $\mu_0=\delta_0$, $\alpha=\mu_n K$, and, conditionally on
$\mathcal F_n$, the new slice is independent Gaussian. The ratio
$W_{n+1}^\beta/W_n^\beta$ is $A_{\beta^2}$ and the next endpoint law is
$\nu_{\beta^2}$. Equation~\eqref{eq:lyapunov-step} therefore says that
\[
 (W_n^\beta)^sH(\mu_n),\qquad n\ge0,
\]
is a nonnegative supermartingale. Every term is integrable by
\eqref{eq:finite-moments} and \eqref{eq:terminal-bounds}. Taking expectations
and using $H\ge1$ yields
\begin{equation}\label{eq:all-N}
 \sup_{\beta^2\in[t_0,t_0+\tau]}\sup_{N\ge0}
 \E[(W_N^\beta)^s]\le H(\delta_0)=C_p^a.
\end{equation}
Taking $\varepsilon=\sqrt{\beta_0^2+\tau}-\beta_0$ proves
Proposition~\ref{prop:stability}.

\subsection{Proof of the moment-spectrum theorem}
\begin{proof}[Proof of Theorem~\ref{thm:main}]
By Proposition~\ref{prop:comparison}, $\ps$ is finite and strictly decreasing
on $(0,\beta_c]$, with limit infinity at zero and value $5/3$ at $\beta_c$.
It remains to prove continuity. At every interior point with $\ps(\beta)>2$,
Proposition~\ref{prop:stability} applies to every real $s\in(2,\ps(\beta))$.
Thus the right limit of the nonincreasing function $\ps$ is at least each such
$s$, and hence equals $\ps(\beta)$. At an interior point with $\ps(\beta)\le2$,
strict decrease and \eqref{eq:critical-value} give $\ps(\beta)>5/3$;
\cite[Theorem~1.2(i)]{Junk} supplies right continuity, including the case
$\ps(\beta)=2$.

Weak disorder holds at every $\beta\in(0,\beta_c]$, so
\eqref{eq:left-continuity} supplies left continuity throughout this domain,
including the critical endpoint. The resulting continuous strictly decreasing
function takes every finite value above $5/3$ by the intermediate value theorem:
for such a value choose a sufficiently small positive $b$ with $\ps(b)$ above
it and use the interval $[b,\beta_c]$. Strict decrease makes the preimage unique.
The value $5/3$ is attained only at $\beta_c$.
\end{proof}

\section*{AI assistance}
OpenAI Codex was used extensively in developing the argument, checking sources
and drafting the text. It also supplied internal reader and referee feedback.
Independent mathematical review is requested.

\begin{thebibliography}{7}
\bibitem{Berger}
Q.~Berger, \emph{Random polymers in random environment: an introduction to
disordered systems}, lecture notes, April 16, 2026.
\url{https://www.math.univ-paris13.fr/~quentin.berger/documents/Disordered-Polymers.pdf}.

\bibitem{ChafaiLehec}
D.~Chafa\"i and J.~Lehec, \emph{Logarithmic Sobolev Inequalities Essentials},
lecture notes, Winter 2023 revision.
\url{https://djalil.chafai.net/docs/M2/chafai-lehec-m2-lsie-lecture-notes.pdf}.

\bibitem{Durrett}
R.~Durrett, \emph{Probability: Theory and Examples}, online version 5,
January 11, 2019.
\url{https://sites.math.duke.edu/~rtd/PTE/PTE5_011119.pdf}.

\bibitem{Junk}
S.~Junk, \emph{Local limit theorem for directed polymers beyond the
$L^2$-phase}, arXiv:2307.05097v3.
\url{https://arxiv.org/html/2307.05097v3}.

\bibitem{Tail}
S.~Junk and H.~Lacoin, \emph{The tail distribution of the partition function
for directed polymers in the weak disorder phase}, arXiv:2405.04335v3.
\url{https://arxiv.org/html/2405.04335v3}.

\bibitem{Critical}
S.~Junk and H.~Lacoin, \emph{Coincidence of critical points for directed
polymers for general environments and random walks}, arXiv:2502.04113v2.
\url{https://arxiv.org/html/2502.04113v2}.

\bibitem{Kemp}
T.~Kemp, \emph{Hypercontractivity in non-commutative holomorphic spaces},
Communications in Mathematical Physics \textbf{259} (2005), 615--637.
\href{https://doi.org/10.1007/s00220-005-1379-5}{doi:10.1007/s00220-005-1379-5}.
Author-generated copy, November~5, 2005:
\url{https://mathweb.ucsd.edu/~tkemp/K-CMP-2005.pdf}.
\end{thebibliography}
\end{document}
